\documentclass[10pt,a4paper]{amsart}
\usepackage{amsmath,amssymb,amsthm,graphicx}
\usepackage{verbatim}
\newenvironment{reference}{\comment}{\endcomment}
%\newenvironment{reference}{}{}
\newenvironment{slogan}{\comment}{\endcomment}
\newenvironment{history}{\comment}{\endcomment}

% For commutative diagrams we use Xy-pic
\usepackage[all]{xy}

% We use 2cell for 2-commutative diagrams.
\xyoption{2cell}
% Theorem environments.
%
\theoremstyle{plain}
\newtheorem{theorem}[subsection]{Theorem}
\newtheorem{proposition}[subsection]{Proposition}
\newtheorem{lemma}[subsection]{Lemma}

\theoremstyle{definition}
\newtheorem{definition}[subsection]{Definition}
\newtheorem{example}[subsection]{Example}
\newtheorem{exercise}[subsection]{Exercise}
\newtheorem{situation}[subsection]{Situation}

\theoremstyle{remark}
\newtheorem{remark}[subsection]{Remark}
\newtheorem{remarks}[subsection]{Remarks}

\numberwithin{equation}{subsection}

% Macros
%
\def\lim{\mathop{\mathrm{lim}}\nolimits}
\def\colim{\mathop{\mathrm{colim}}\nolimits}
\def\Spec{\mathop{\mathrm{Spec}}}
\def\Hom{\mathop{\mathrm{Hom}}\nolimits}
\def\Ext{\mathop{\mathrm{Ext}}\nolimits}
\def\SheafHom{\mathop{\mathcal{H}\!\mathit{om}}\nolimits}
\def\SheafExt{\mathop{\mathcal{E}\!\mathit{xt}}\nolimits}
\def\Sch{\mathit{Sch}}
\def\Mor{\mathop{\mathrm{Mor}}\nolimits}
\def\Ob{\mathop{\mathrm{Ob}}\nolimits}
\def\Sh{\mathop{\mathit{Sh}}\nolimits}
\def\NL{\mathop{N\!L}\nolimits}
\def\CH{\mathop{\mathrm{CH}}\nolimits}
\def\proetale{{pro\text{-}\acute{e}tale}}
\def\etale{{\acute{e}tale}}
\def\QCoh{\mathit{QCoh}}
\def\Ker{\mathop{\mathrm{Ker}}}
\def\Im{\mathop{\mathrm{Im}}}
\def\Coker{\mathop{\mathrm{Coker}}}
\def\Coim{\mathop{\mathrm{Coim}}}

% Boxtimes
%
\DeclareMathSymbol{\boxtimes}{\mathbin}{AMSa}{"02}

%
% Macros for moduli stacks/spaces
%
\def\QCohstack{\mathcal{QC}\!\mathit{oh}}
\def\Cohstack{\mathcal{C}\!\mathit{oh}}
\def\Spacesstack{\mathcal{S}\!\mathit{paces}}
\def\Quotfunctor{\mathrm{Quot}}
\def\Hilbfunctor{\mathrm{Hilb}}
\def\Curvesstack{\mathcal{C}\!\mathit{urves}}
\def\Polarizedstack{\mathcal{P}\!\mathit{olarized}}
\def\Complexesstack{\mathcal{C}\!\mathit{omplexes}}
% \Pic is the operator that assigns to X its picard group, usage \Pic(X)
% \Picardstack_{X/B} denotes the Picard stack of X over B
% \Picardfunctor_{X/B} denotes the Picard functor of X over B
\def\Pic{\mathop{\mathrm{Pic}}\nolimits}
\def\Picardstack{\mathcal{P}\!\mathit{ic}}
\def\Picardfunctor{\mathrm{Pic}}
\def\Deformationcategory{\mathcal{D}\!\mathit{ef}}
\usepackage{cleveref}
\crefname{originaldefinition}{original definition}{original definitions}
\Crefname{originaldefinition}{Original definition}{Original definitions}
\makeatletter
\newlabel{definition-scheme-theoretic-image}{{082Y}{source}{}{sourcepointer-1}{}}
\newlabel{definition-scheme-theoretic-image@cref}{{[originaldefinition][0][]082Y}{[1][0][]source}{}{}{}}
\makeatother
\crefname{originallemma}{original lemma}{original lemmas}
\Crefname{originallemma}{Original lemma}{Original lemmas}
\makeatletter
\newlabel{lemma-image-proper-is-proper}{{08AJ}{source}{}{sourcepointer-2}{}}
\newlabel{lemma-image-proper-is-proper@cref}{{[originallemma][0][]08AJ}{[1][0][]source}{}{}{}}
\makeatother
\crefname{originallemma}{original lemma}{original lemmas}
\Crefname{originallemma}{Original lemma}{Original lemmas}
\makeatletter
\newlabel{lemma-integral-universally-closed}{{0415}{source}{}{sourcepointer-3}{}}
\newlabel{lemma-integral-universally-closed@cref}{{[originallemma][0][]0415}{[1][0][]source}{}{}{}}
\makeatother
\crefname{originallemma}{original lemma}{original lemmas}
\Crefname{originallemma}{Original lemma}{Original lemmas}
\makeatletter
\newlabel{lemma-quasi-compact-permanence}{{03KS}{source}{}{sourcepointer-4}{}}
\newlabel{lemma-quasi-compact-permanence@cref}{{[originallemma][0][]03KS}{[1][0][]source}{}{}{}}
\makeatother
\crefname{originallemma}{original lemma}{original lemmas}
\Crefname{originallemma}{Original lemma}{Original lemmas}
\makeatletter
\newlabel{lemma-quasi-compact-scheme-theoretic-image}{{082Z}{source}{}{sourcepointer-5}{}}
\newlabel{lemma-quasi-compact-scheme-theoretic-image@cref}{{[originallemma][0][]082Z}{[1][0][]source}{}{}{}}
\makeatother
\crefname{originallemma}{original lemma}{original lemmas}
\Crefname{originallemma}{Original lemma}{Original lemmas}
\makeatletter
\newlabel{lemma-universally-closed-permanence}{{04NX}{source}{}{sourcepointer-6}{}}
\newlabel{lemma-universally-closed-permanence@cref}{{[originallemma][0][]04NX}{[1][0][]source}{}{}{}}
\makeatother
\crefname{originalremark}{original remark}{original remarks}
\Crefname{originalremark}{Original remark}{Original remarks}
\makeatletter
\newlabel{remark-factorization-quasi-compact-quasi-separated}{{081X}{source}{}{sourcepointer-7}{}}
\newlabel{remark-factorization-quasi-compact-quasi-separated@cref}{{[originalremark][0][]081X}{[1][0][]source}{}{}{}}
\makeatother
\title{Pushforward of the structure sheaf along a universally closed morphism is integral over the base}
\author{The Stacks Project Authors}
\date{stacks-project:spaces-morphisms.tex:local-snapshot-20261012; GNU FDL 1.2 or later}
\begin{document}
\maketitle
\setcounter{subsection}{1}
\noindent\textit{Editorial note.} All delivered mathematical text is the author's own, copied from the pinned Stacks Project chapter file spaces-morphisms.tex; no statement or proof step was rewritten, reordered or removed.\par\medskip
\noindent\textit{Editorial note.} The proof is the author's reduction to a single section h and produces the monic equation explicitly; the properness and integrality inputs are applied inside the counted span.\par\medskip
\noindent\textit{Editorial note.} The seven references of the counted spans are recorded as non-load-bearing navigation pointers; the argument that consumes them stays inside the counted proof span.\par\medskip
\noindent\textit{Editorial note.} Preamble: the three native excerpt spans of sources/stacks-spaces-morphisms/preamble.tex (lines 8-13, 14-18 and 154-217). The chapter file itself carries no preamble text.\par\medskip
\noindent\textit{Allowed standard prerequisites (editorial).}\begin{itemize}
\item The factorization of a quasi-compact quasi-separated morphism through the relative spectrum of its pushforward. Remark remark-factorization-quasi-compact-quasi-separated, used at proof lines 10150-10155. Fixes the factorization that the lemma identifies with the normalization.
\item Scheme theoretic image and its density for quasi-compact morphisms. Definition definition-scheme-theoretic-image and Lemma lemma-quasi-compact-scheme-theoretic-image, used at proof lines 10169-10175. Produces Z and the density of X -> Z.
\item Permanence of universal closedness through a factorization. Lemma lemma-universally-closed-permanence, used at proof lines 10176-10177. Makes X -> Z closed, so that h(X) = Z.
\item Images of proper morphisms are proper, and an affine universally closed morphism is integral. Lemmas lemma-image-proper-is-proper and lemma-integral-universally-closed, used at proof lines 10178-10182. Makes Z -> Y integral over an affine base.
\item Permanence of quasi-compactness for the morphism h. Lemma lemma-quasi-compact-permanence, used at proof lines 10171-10173. Shows h is quasi-compact so that the scheme theoretic image argument applies.
\end{itemize}\medskip
\ifdefined\hypertarget\hypertarget{sourcepointer-1}{}\fi
\noindent\textit{Source reference.} Original definition 082Y: focused-seeds/AGP/sources/stacks-snapshot-20261011/tags.txt, source lines 10268--10268. Navigation pointer into the Stacks Project; the standard fact it supplies is applied inside the counted proof span and the argument resting on it stays there.\par\medskip
\ifdefined\hypertarget\hypertarget{sourcepointer-2}{}\fi
\noindent\textit{Source reference.} Original lemma 08AJ: focused-seeds/AGP/sources/stacks-snapshot-20261011/tags.txt, source lines 10534--10534. Navigation pointer into the Stacks Project; the standard fact it supplies is applied inside the counted proof span and the argument resting on it stays there.\par\medskip
\ifdefined\hypertarget\hypertarget{sourcepointer-3}{}\fi
\noindent\textit{Source reference.} Original lemma 0415: focused-seeds/AGP/sources/stacks-snapshot-20261011/tags.txt, source lines 5280--5280. Navigation pointer into the Stacks Project; the standard fact it supplies is applied inside the counted proof span and the argument resting on it stays there.\par\medskip
\ifdefined\hypertarget\hypertarget{sourcepointer-4}{}\fi
\noindent\textit{Source reference.} Original lemma 03KS: focused-seeds/AGP/sources/stacks-snapshot-20261011/tags.txt, source lines 4688--4688. Navigation pointer into the Stacks Project; the standard fact it supplies is applied inside the counted proof span and the argument resting on it stays there.\par\medskip
\ifdefined\hypertarget\hypertarget{sourcepointer-5}{}\fi
\noindent\textit{Source reference.} Original lemma 082Z: focused-seeds/AGP/sources/stacks-snapshot-20261011/tags.txt, source lines 10269--10269. Navigation pointer into the Stacks Project; the standard fact it supplies is applied inside the counted proof span and the argument resting on it stays there.\par\medskip
\ifdefined\hypertarget\hypertarget{sourcepointer-6}{}\fi
\noindent\textit{Source reference.} Original lemma 04NX: focused-seeds/AGP/sources/stacks-snapshot-20261011/tags.txt, source lines 6095--6095. Navigation pointer into the Stacks Project; the standard fact it supplies is applied inside the counted proof span and the argument resting on it stays there.\par\medskip
\ifdefined\hypertarget\hypertarget{sourcepointer-7}{}\fi
\noindent\textit{Source reference.} Original remark 081X: focused-seeds/AGP/sources/stacks-snapshot-20261011/tags.txt, source lines 10232--10232. Navigation pointer into the Stacks Project; the standard fact it supplies is applied inside the counted proof span and the argument resting on it stays there.\par\medskip
\begin{lemma}
\label{lemma-normalization-in-universally-closed}
Let $S$ be a scheme.
Let $f : X \to Y$ be a quasi-compact, quasi-separated and
universally closed morphisms of algebraic spaces over $S$.
Then $f_*\mathcal{O}_X$ is integral over $\mathcal{O}_Y$. In other
words, the normalization of $Y$ in $X$ is equal to the factorization
$$
X \longrightarrow \underline{\Spec}_Y(f_*\mathcal{O}_X)
\longrightarrow Y
$$
of Remark \ref{remark-factorization-quasi-compact-quasi-separated}.
\end{lemma}

\begin{proof}
The question is \'etale local on $Y$, hence we may reduce to the
case where $Y = \Spec(R)$ is affine. Let $h \in \Gamma(X, \mathcal{O}_X)$.
We have to show that $h$ satisfies a monic equation over $R$. Think of $h$
as a morphism as in the following commutative diagram
$$
\xymatrix{
X \ar[rr]_h \ar[rd]_f & & \mathbf{A}^1_Y \ar[ld] \\
& Y &
}
$$
Let $Z \subset \mathbf{A}^1_Y$ be the scheme theoretic image of $h$,
see Definition \ref{definition-scheme-theoretic-image}.
The morphism $h$ is quasi-compact as $f$ is quasi-compact and
$\mathbf{A}^1_Y \to Y$ is separated, see
Lemma \ref{lemma-quasi-compact-permanence}.
By Lemma \ref{lemma-quasi-compact-scheme-theoretic-image} the
morphism $X \to Z$ has dense image on underlying topological spaces. By
Lemma \ref{lemma-universally-closed-permanence} the morphism
$X \to Z$ is closed. Hence $h(X) = Z$ (set theoretically).
Thus we can use
Lemma \ref{lemma-image-proper-is-proper}
to conclude that $Z \to Y$ is universally closed (and even proper).
Since $Z \subset \mathbf{A}^1_Y$, we see that $Z \to Y$ is affine
and proper, hence integral by Lemma \ref{lemma-integral-universally-closed}.
Writing $\mathbf{A}^1_Y = \Spec(R[T])$ we conclude that
the ideal $I \subset R[T]$ of $Z$ contains a monic polynomial
$P(T) \in R[T]$. Hence $P(h) = 0$ and we win.
\end{proof}

\end{document}
